\ifx\allfiles\undefined
\documentclass[12pt, a4paper, oneside, UTF8]{ctexbook}
\def\path{../config}
\input{\path/package.tex}

% 定理环境设置
\input{\path/theorem1_zh.tex}

\input{\path/custom.tex}

% 修改参数改变封面样式，0 默认原始封面、内置其他1、2、3种封面样式
\def\myIndex{3}


\ifnum\myIndex>0
    \input{\path/cover_package_\myIndex} 
\fi

\def\myTitle{高等数学笔记}
\def\myAuthor{M.K.}
\def\myDateCover{\today}
\def\myDateForeword{\today}
\def\myForeword{前言}
\def\myForewordText{
    
    这是一个基于\LaTeX{}模板的高数笔记．
    内容参考了包括但不限于以下资料：
    \begin{itemize}
        \item 《高等数学》（同济大学出版社）第七版
        \item 《吉米多维奇高等数学习题精选精讲》（山东科学技术出版社）第二版
        \item 《高等数学作业集》（哈尔滨工业大学出版社）第一版
        \item \href{https://space.bilibili.com/1035929235}{B站UP主：一高数}的高数课程
        \item \href{https://space.bilibili.com/42428180}{B站UP主：考研竞赛数学}的试题讲解与高数课程
        \item \href{https://space.bilibili.com/391930545}{B站UP主：Maki的完美算数教室}的高数内容

    \end{itemize}

    封面来源于\href{https://www.pixiv.net/users/3952}{おにねこ}的作品\href{https://www.pixiv.net/artworks/136551599}{雨雫とプレアデス}
    如有侵权请联系我删除．
}
\def\mySubheading{副标题}


\begin{document}
% \input{../config/cover}
\else
\fi
\chapter{二重极限}
\section{二重极限不存在的证明}

\begin{theorem}[单路径法]
若存在一条路径 $L$ 使得点 $P(x,y)$ 沿 $L$ 趋于 $P_0(x_0,y_0)$ 时，函数 $f(x,y)$ 的极限不存在，则极限 $\lim\limits_{(x,y) \to (x_0,y_0)} f(x,y)$ 不存在．
\end{theorem}


\begin{theorem}[双路径法]
若存在两条不同的路径 $L_1, L_2$ 使得点 $P(x,y)$ 分别沿这两条路径趋于 $P_0(x_0,y_0)$ 时，函数 $f(x,y)$ 的极限值不相等，即
\[ \lim_{\substack{(x,y) \to (x_0,y_0) \\ (x,y) \in L_1}} f(x,y) \neq \lim_{\substack{(x,y) \to (x_0,y_0) \\ (x,y) \in L_2}} f(x,y) \]
则极限 $\lim\limits_{(x,y) \to (x_0,y_0)} f(x,y)$ 不存在．
\end{theorem}


\begin{theorem}[$y=kx^p$路径法]
对于 $(x,y) \to (0,0)$ 的情形，选取路径 $y=kx^p$（$p$ 为常数），若极限
\[ \lim_{\substack{(x,y) \to (0,0) \\ y=kx^p}} f(x,y) = \varphi(k) \]
依赖于参数 $k$，则极限 $\lim\limits_{(x,y) \to (0,0)} f(x,y)$ 不存在．
\end{theorem}

\begin{corollary}
    对于形如 $\lim\limits_{(x,y) \to (0,0)} \dfrac{x^p y^q}{x^m+y^n}\;(m, n\in \mathbb{Z})^+,\,p,q>0$ 有下列结论：
    \begin{enumerate}[label=(\arabic*)]
        \item 若$m,n$中有奇数，极限不存在；
        \item 若$m,n$均为偶数，则：\begin{enumerate}[label=\roman*.]
            \item 若$\dfrac{p}{m}+\dfrac{q}{n}>1$，极限为0；
            \item 若$\dfrac{p}{m}+\dfrac{q}{n}\leq 1$，选取路径$y=kx^{\frac{m-p}{q}}$，极限不存在．
            \end{enumerate}
    \end{enumerate}
\end{corollary}


\begin{theorem}[极坐标法]
作极坐标变换 $\begin{cases} x=x_0+r\cos\theta \\ y=y_0+r\sin\theta \end{cases}$，若极限
\[ \lim_{r \to 0^+} f(x_0+r\cos\theta, y_0+r\sin\theta) = \psi(\theta) \]
的值依赖于 $\theta$，或者该极限不存在，则极限 $\lim\limits_{(x,y) \to (x_0,y_0)} f(x,y)$ 不存在．
\end{theorem}
\begin{example}
    证明下列极限不存在：
    \begin{tasks}[label=(\arabic*)](2)
        \task $\lim\limits_{(x,y) \to (0,0)} \dfrac{x+y}{xy}$；
        \task $\lim\limits_{(x,y) \to (0,0)} \dfrac{xy}{x^2+y^2}$；
        \task $\lim\limits_{(x,y) \to (0,0)} \dfrac{xy^2}{x^2+y^4}$；
        \task $\lim\limits_{(x,y) \to (0,0)} \dfrac{x+y}{x-y}$．
    \end{tasks}
\end{example}
\begin{proof}
    \begin{enumerate}[label=(\arabic*)]
        \item 选取路径 $y=x$，则 $\lim\limits_{\substack{(x,y) \to (0,0) \\ y=x}} \dfrac{x+y}{xy} = \lim\limits_{x \to 0} \dfrac{2x}{x^2} = \lim\limits_{x \to 0} \dfrac{2}{x}$ 不存在，故原极限不存在．
        \item 法一：选取路径 $y=x$和 $y=-x$，则
        \[ \lim_{\substack{(x,y) \to (0,0) \\ y=x}} \dfrac{xy}{x^2+y^2} = \lim_{x \to 0} \dfrac{x^2}{2x^2} = \dfrac{1}{2}, \]
        \[ \lim_{\substack{(x,y) \to (0,0) \\ y=-x}} \dfrac{xy}{x^2+y^2} = \lim_{x \to 0} \dfrac{-x^2}{2x^2} = -\dfrac{1}{2}, \]
        故原极限不存在．\\
        法二：选取路径 $y=kx$，则
        \[ \lim_{\substack{(x,y) \to (0,0) \\ y=kx}} \dfrac{xy}{x^2+y^2} = \lim_{x \to 0} \dfrac{kx^2}{x^2+k^2x^2} = \dfrac{k}{1+k^2}, \]
        该极限依赖于参数 $k$，故原极限不存在．\\
        法三：作极坐标变换 $x=r\cos\theta, y=r\sin\theta$，则
        \[ \lim_{r \to 0^+} \dfrac{(r\cos\theta)(r\sin\theta)}{(r\cos\theta)^2+(r\sin\theta)^2} = \lim_{r \to 0^+} \dfrac{r^2\cos\theta\sin\theta}{r^2} = \cos\theta\sin\theta, \]
        该极限依赖于参数 $\theta$，故原极限不存在．
        \item 选取路径 $x = ky^2$，则
        \[ \lim_{\substack{(x,y) \to (0,0) \\ x=ky^2}} \dfrac{xy^2}{x^2+y^4} = \lim_{y \to 0} \dfrac{(ky^2)y^2}{(ky^2)^2+y^4} = \lim_{y \to 0} \dfrac{ky^4}{k^2y^4+y^4} = \dfrac{k}{k^2+1}, \]
        该极限依赖于参数 $k$，故原极限不存在．
        \item 选取路径 $y = x + x^2$，则
        \[ \lim_{\substack{(x,y) \to (0,0) \\ y=x+x^2}} \dfrac{x+y}{x-y} = \lim_{x \to 0} \dfrac{x+(x+x^2)}{x-(x+x^2)} = \lim_{x \to 0} \dfrac{2x+x^2}{-x^2} =\infty \]
        不存在，故原极限不存在．
    \end{enumerate}
\end{proof}
\section{二重极限的计算}
\begin{theorem}
    二重极限常有下列方法：
    \begin{enumerate}[label=(\arabic*)]
        \item 连续$\Rightarrow$直接代入法；
        \item 等价无穷小替换法；
        \item 夹逼准则；
        \item $0\times$有界量法；
        \item 极坐标变换法．
    \end{enumerate}
\end{theorem}
\begin{example}
    计算下列极限：
    \begin{tasks}[label=(\arabic*)](3)
        \task $\lim\limits_{(x,y) \to (0,0)} \dfrac{xy}{\sqrt{x^2+y^2}}$；
        \task $\lim\limits_{(x,y) \to (0,0)} \dfrac{\ln(1+x^2y^4)}{x^2+y^2}$；
        \task $\lim\limits_{(x,y) \to (0,0)} \dfrac{x^2 +y^2}{|x| + |y|}$；
    \end{tasks}
\end{example}
\begin{proof}
    \begin{enumerate}[label=(\arabic*)]
        \item \[0<=\left| \dfrac{xy}{\sqrt{x^2+y^2}} \right| \leq |x| \to 0 (x,y) \to (0,0)\]
        故原极限为0．
        \item \begin{align*}
            \lim\limits_{(x,y) \to (0,0)} \dfrac{\ln(1+x^2y^4)}{x^2+y^2} & = \lim\limits_{(x,y) \to (0,0)} \dfrac{x^2y^4}{x^2+y^2} \quad (\text{因为} \ln(1+u) \sim u, u \to 0) \\
            & = \lim\limits_{(x,y) \to (0,0)} y^4 \cdot \dfrac{x^2}{x^2+y^2} 
        \end{align*}
        其中 $0 < \dfrac{x^2}{x^2+y^2} \leq 1$，且 $\lim\limits_{(x,y) \to (0,0)} y^4 = 0$，由 $0 \times$ 有界量法，原极限为0．\\
        \item 法一：
        \[ \lim\limits_{(x,y) \to (0,0)} \dfrac{x^2 +y^2}{|x| + |y|}= \lim\limits_{(x,y) \to (0,0)} |x| \cdot \dfrac{|x|}{|x| + |y|} + |y| \cdot \dfrac{|y|}{|x| + |y|} \]
        其中 $0 \leq \dfrac{|x|}{|x| + |y|} \leq 1$，$0 \leq \dfrac{|y|}{|x| + |y|} \leq 1$，且 $\lim\limits_{(x,y) \to (0,0)} |x| = 0$，$\lim\limits_{(x,y) \to (0,0)} |y| = 0$，由 $0 \times$ 有界量法，原极限为0．\\
        法二：作极坐标变换 $x=r\cos\theta, y=r\sin\theta$，则
        \[ \lim_{r \to 0^+} \dfrac{(r\cos\theta)^2 +(r\sin\theta)^2}{|r\cos\theta| + |r\sin\theta|} = \lim_{r \to 0^+} \dfrac{r^2}{r(|\cos\theta| + |\sin\theta|)} = \lim_{r \to 0^+} \dfrac{r}{|\cos\theta| + |\sin\theta|} = 0. \]
        故原极限为0．
    \end{enumerate}
\end{proof}
\begin{example}
    求二元函数 $f(x,y) = \begin{cases}
        \dfrac{2xy^2}{x^2+y^2}-\dfrac{2x^3y^2}{(x^2+y^2)^2}, & (x,y) \neq (0,0) \\
        0, & (x,y) = (0,0)
    \end{cases}$ 在点 $(0,0)$ 处是否连续．
\end{example}
\begin{solution}
    \begin{align*}
        \lim\limits_{(x,y) \to (0,0)} f(x,y) & = \lim\limits_{(x,y) \to (0,0)} \left( \dfrac{2xy^2}{x^2+y^2}-\dfrac{2x^3y^2}{(x^2+y^2)^2} \right) \\
        & = \lim\limits_{(x,y) \to (0,0)} 2x\cdot\dfrac{y^2}{x^2+y^2}-x\dfrac{2x^2y^2}{(x^2+y^2)^2}
    \end{align*}
    其中 $0 \leq \dfrac{y^2}{x^2+y^2} \leq 1$，$0 \leq \dfrac{2x^2y^2}{(x^2+y^2)^2} \leq 1$，且 $\lim\limits_{(x,y) \to (0,0)} 2x = 0$，$\lim\limits_{(x,y) \to (0,0)} x = 0$，由 $0 \times$ 有界量法，原极限为0．\\
    因为 $\lim\limits_{(x,y) \to (0,0)} f(x,y) = f(0,0) = 0$，故函数 $f(x,y)$ 在点 $(0,0)$ 处连续．
\end{solution}
\subsection{二次极限}
\begin{example}
    求极限 $\lim\limits_{x \to 0}\lim\limits_{y \to 0} (x\sin\dfrac{1}{y}+y\sin\dfrac{1}{x})$
\end{example}
\begin{solution}
    由于$\lim\limits_{y \to 0} x\sin\dfrac{1}{y}$ 不存在，故整个极限不存在．
\end{solution}
\ifx\allfiles\undefined
\end{document}
\fi